\section{Introduction}

% The figure of 150.6 km² represents the raw total land area of the entire municipality,
% not just the urbanized footprint. Because of this, the average figure of roughly 1,385 inhabitants/km² is indeed a massive underestimation of the actual population density
% in the areas where people actually live.
% Here is how the territorial reality and the legal framework explain this discrepancy:
% 1. The Reality of the Terrain
% A vast portion of the municipality of Santa Cruz de Tenerife is composed of the Macizo de Anaga. This area, known as the Parque Rural de Anaga, is a declared Biosphere Reserve and features steep, mountainous terrain that is largely uninhabitable. Because this massive natural expanse is legally part of the municipality's total surface area, dividing the total population by the total km² artificially dilutes the density metric. If you were to look exclusively at the urbanized footprint (the city center, the coastal strip, and the southwest neighborhoods), the population density would be drastically higher, resembling that of major dense European cities.   2. Legal Protection and Land ClassificationThe reason this steep mountainous area remains uninhabited is not just topographical, but strictly regulated by law. If you reference the file named "Ley del Suelo, Planes de Ordenación Urbana", you will see that territorial and urban planning instruments mandate the strict delimitation and protection of these spaces.According to the law, island and municipal planning must identify and classify specific zones of the territory that are to be excluded from development, including:Areas that must be preserved from the urbanization and building process because transforming them would be incompatible with sustainable development.   Areas that must be preserved due to their existing or potential landscape, agricultural, or heritage value.   Areas that must be excluded from the urbanization process to prevent risks—such as geological, seismic, or meteorological risks (which is highly relevant for steep ravines and mountainsides).   Therefore, under the "Ley del Suelo, Planes de Ordenación Urbana", the mountain ranges of Anaga are strictly classified as protected rural land (suelo rústico / espacios naturales protegidos), effectively fencing them off from urban expansion.

The figure of 150.6~km\textsuperscript{2} represents the raw total land area of the entire municipality of Santa Cruz de Tenerife, not merely the urbanized footprint. For this reason, the municipal average indicator of approximately 1,385~inhabitants/km\textsuperscript{2} constitutes a substantial underestimation of the actual population density sustained in the sectors where people actually live.

This marked discrepancy between the aggregate statistical figure and the physical spatial reality is explained by two fundamental structural conditions:

\begin{enumerate}
  \item \textbf{The physical and orographic reality of the terrain:} A vast expanse of the municipality is occupied by the Anaga Massif (\emph{Macizo de Anaga}). This area, protected under the status of Rural Park and designated a UNESCO Biosphere Reserve, features steep, mountainous topography with severe slopes that is materially uninhabitable and hostile to urban development. Because this massive natural expanse is legally included within the total surface area of the municipality, the arithmetic calculation dividing total population by municipal land area artificially dilutes the density index. If the analysis is restricted to the actual urbanized footprint (the historic center and extensions, the coastal strip, and the southwest neighborhoods), net population density increases dramatically, reaching figures comparable to the densest European metropolises.
  \item \textbf{Legal protection and land classification:} The absence of urban development across mountainous areas is not merely due to topographical barriers, but to strict legal safeguards established under land legislation and territorial and urban planning instruments. Canary Islands legislation mandates that spatial plans identify, delimit, and strictly protect specific zones excluded from development, including:
  \begin{itemize}
    \item Land that must be preserved from urbanization and construction because transformation would be incompatible with sustainable territorial development.
    \item Areas that must be preserved due to their intrinsic landscape, agricultural, forestry, ecological, natural, or heritage values.
    \item Land that must be excluded from urbanization specifically to prevent and mitigate severe natural and geotechnical risks (geological, seismic, or meteorological hazards, which are especially acute along steep slopes and ravine beds).
  \end{itemize}
  Under this urban planning regime, the ridges and foothills of the Anaga Massif are strictly classified as Protected Rural Land (\emph{Suelo Rústico Protegido}, integrated into the Canary Network of Protected Natural Spaces), establishing an impenetrable geographic and legal barrier against municipal urban expansion.
\end{enumerate}

Under these structural constraints of physical containment and legal limitation of habitable space, the following sections examine the demographic, socioeconomic, and regulatory forces driving urban planning and spatial organization in Santa Cruz de Tenerife from 2011 to the present.

\subsection{Tourism and Perception: The Disparity Between Lodging and Visitor Influx}

% 1. The 2011–2012 Population Drop: Factual Exodus or Statistical Artifact?
% It was purely an administrative and statistical artifact, not a physical exodus of 15,000 residents.

% What actually happened:
% The 2011 Census Audit: In 2011–2012, the INE conducted Spain's decennial Censo de Población y Viviendas 2011, transitioning from door-to-door paper forms to cross-referencing administrative databases (Padrón, Seguridad Social, and the Ministry of the Interior's National DNI/NIE database).
% The "Habitantes Fantasma" of Santa Cruz: When the INE audited Santa Cruz de Tenerife's municipal register, it discovered an accumulated backlog of invalid, duplicated, or non-existent DNI entries that had been carried over since 1996.
% The Official Purge: The INE rejected the city's official submission and forced an administrative cleanup of 15,306 entries on 1 January 2012 (a sudden drop from 222,271 down to 206,965, or -6.89%).
% The Political & Fiscal Scandal: Because municipal funding allocations (such as the Fondo Canario de Financiación Municipal, state tax shares, and Canary REF revenue) were distributed per capita based on padrón figures, this purge caused a major municipal and political controversy across the Canary Islands.

% 2. Would plotting these series showcase the need for efficient land use management?
% Yes, absolutely. In fact, it provides a much more sophisticated urban planning argument than the typical "runaway demographic growth" narrative.

% Here is why and how it builds your thesis:

% A. The Paradox of "Flat Growth" vs. "Severe Urban Tension"
% Looking at the data extracted from your files:

% Resident Population (population_all.json): Post-2012, population has been virtually flat (from 206,965 in 2012 to 211,498 in 2025, barely +2.18% over 13 years).
% Touristic Lodging Population (alojatur_poblacion_tur_etc.json): Santa Cruz is the provincial administrative and port capital, not a resort enclave like Arona or Adeje. Its regulated tourist population is modest (oscillating between ~900 and ~1,700 daily equivalent tourists), though pressure from vivienda vacacional (VV) and daytime excursionists in central quarters has surged recently.
% Household Income (renta.json): Average gross income rose steadily from 26,124 € (2013) to 28,662 € (2021).
% B. The Planning Thesis (The Supply Bottleneck)
% In conventional planning debates, housing shortages and spatial stress are blamed on explosive demographic booms. In Santa Cruz, the data shows the opposite:

% Shrinking Household Size: Even with flat population, households have unbundled (more single-person and two-person homes), meaning the city requires significantly more housing units just to house the same number of people.
% Economic & Commercial Vitality: Rising income drives demand for better housing quality, logistics space, and tertiary services.
% The Artificial Bottleneck (Planning Paralysis): Because the PGO 2013 was annulled by the TSJC (Sentence 116/2014) and the 2022 provisional ordinance was suspended (Sentence 70/2023), the city is legally anchored to the obsolete 2005 plan.
% Land transformation in Sectorized Urbanizable Land (SUSO) and Unconsolidated Urban Land (SUNC) is virtually frozen.
% Private and public developers cannot deliver new housing stock or complete strategic urban stitching (e.g., around the Barranco de Santos or Cabo Llanos).

Since the beginning of available data (2011--2025), the population of Santa Cruz de Tenerife has remained stable between 205,000 and 210,000 inhabitants.
Regulated tourist lodging population recorded by the Canary Institute of Statistics (ISTAC) barely reaches 2,000 equivalent residents.
Although it has grown by 50\% over the period analyzed, rapid growth of an initially tiny figure still results in a small absolute number.

\begin{figure}[htbp]
  \centering
  \includegraphics[width=0.95\linewidth]{figures/demographics_trend.pdf}
  \caption{Evolution of Resident and Tourist Population in Santa Cruz de Tenerife (2011--2025). Source: Author's elaboration based on ISTAC and INE data.}
  \label{fig:demografia_trend}
\end{figure}

%  there is an enormous disparity between the perceived tourist presence in Santa Cruz de Tenerife and the actual resident lodging demand, driven almost entirely by day-visitors (cruceristas and excursionists).

% Here is the breakdown based on official data from the Autoridad Portuaria de Santa Cruz de Tenerife (Puertos de Tenerife), Aena, and the ISTAC:

% 1. The Numbers: Cruise Passengers vs. Air Tourists
% Category  Volume (Annual)  Modality / Behavior in Santa Cruz  Impact on Housing Stock
% Cruise Tourists (Puerto de S/C)  ~972.000 (2024)
% (817.647 en 2023)  Day-trippers (Transeúntes de escala):
% • ~65% (~630.000) disembark for 4–8 hours in the city center and sleep on the ship.
% • ~35% (~340.000) are puerto base (start/end itinerary).  Zero housing impact (they sleep on the ship; they only consume sidewalks, terraces, and retail).
% Air Tourists (Tenerife Island Total)  ~6.500.000 (2023–2024)
% (TFS + TFN)  Concentrated in resort enclaves:
% • 75%–80% stay in the South (Adeje, Arona, Guía de Isora).
% • 15% in the North (Puerto de la Cruz).
% • Barely 2%–3% stay overnight in Santa Cruz.  Localized pressure in southern municipalities, not Santa Cruz.
% Regulated Overnight Tourists in Santa Cruz  ~1.500 – 1.700
% (Población equivalente diaria, Alojatur)  Urban / business / cultural hotel lodging with an average stay of 2–3 nights.  Marginal (<0.8% of the municipal resident population).
% Day Excursionists from the South/North  Hundreds of thousands  Tourists staying in Playa de las Américas, Adeje, or Puerto de la Cruz visiting Santa Cruz for shopping, museums, or Anaga.  Zero housing impact (they return to their southern hotels before dinner).
% 2. Why Perception Disconnects from Reality
% Spatial Concentration in the Historic Core: Cruise ships dock directly at the Muelle Sur and Muelle de Ribera, just steps from Plaza de España, Calle Castillo, Plaza de la Candelaria, and the Mercado de Nuestra Señora de África.
% On peak winter days (October–April), when 2 or 3 large liners dock simultaneously, 8,000 to 12,000 passengers flood a narrow, pedestrianized central corridor between 09:00 and 17:00.
% This creates a vivid visual impression of intense "over-tourism", even though not a single one of these visitors rents an apartment or stays in a local hotel bed.
% Daytime Excursionists: Tourists vacationing in the South take rental cars or TITSA express buses (Línea 110/111) to spend the day in the capital, shopping or touring the Anaga Biosphere Reserve, concentrating human traffic in central streets without registering in municipal lodging statistics.
% 3. Urban Planning & Housing Takeaway (For your Introduction)
% This distinction sharpens your argument in

% sections/15_introduccion.tex
% :

% The Fallacy: Attributing the housing shortage and spatial tension in Santa Cruz de Tenerife to "uncontrolled tourist lodging."
% The Reality:
% Nearly 1 million cruise passengers per year create high pedestrian pressure on public space, street cleaning, and retail, but do not compete for residential housing.
% Over-night lodging in the capital remains tiny (~1,600 equivalent residents).
% Therefore, the housing affordability crisis is almost entirely an internal supply crisis:
% Demographic unbundling: Flat population, but smaller households requiring more individual units.
% Regulatory paralysis: An annulled PGO (2013) and suspended provisional ordinance (2022) that have frozen the legal transformation of new residential sectors (SUSO) and urban regeneration units (SUNC).

In sharp contrast to this low regulated lodging pressure, there is a strong civic perception of tourist saturation.
This phenomenon is driven primarily by day tourism:
the Port of Santa Cruz de Tenerife exceeds 900,000 cruise passengers annually (with more than 600,000 port-of-call transit passengers disembarking for several hours in the historic center),
joined by hundreds of thousands of day excursionists arriving from the southern resort enclaves.
This floating daytime influx intensifies the use of public spaces and urban services during the day, but consumes zero permanent residential housing stock,
confirming that the acute housing shortage in the municipality is not caused by mass tourist lodging absorption.

\subsection{Immigration and Demographic Flows}

% 1. The Latin American Inversion (Historical Flow to America Reversed)
% Status: 100% Factual and Accurate.
% The Historical Context: For over a century, the Canary Islands were an emigration society. Tens of thousands of Canarians emigrated to America—chiefly Cuba (early 20th century) and Venezuela (1950s–1980s, historically known as "la octava isla").
% The Modern Reversal:
% In the last 15 years, that flow has completely inverted. Due to political and economic instability in Latin America (notably Venezuela, Colombia, and Cuba), hundreds of thousands of people have migrated to the Canary Islands.
% A substantial portion are returnees and descendants (children and grandchildren of Canarians) who already hold or easily acquire Spanish nationality under the Ley de Memoria Democrática (Ley de Nietos).
% In Santa Cruz de Tenerife, registered foreign nationals alone rose from 4,473 in 2000 to a record 16,776 in 2025 (ISTAC), not including the tens of thousands who hold Spanish passports by descent.
% Impact on Housing Demand: Very High and Direct.
% These individuals and families settle permanently as long-term residents (vivienda habitual).
% They enter the formal local labor market (healthcare, services, commerce, education) and compete directly in the residential rental and purchase market across Santa Cruz (Ofra, Añaza, La Salud, Tome Cano, Centro).
% 2. Irregular Migration from Africa (Cayucos / Pateras)
% Status: Crucial Correction — It does NOT significantly impact residential housing demand.
% Why it does not affect the private housing market:
% A Transit Phenomenon, Not Settlement: The Ruta Canaria is an emergency transit border route. Over 90% to 95% of adult irregular migrants arriving by sea do not settle in the Canary Islands. The Ministry of Inclusion systematically organizes humanitarian transfers (derivaciones) to mainland Spain (la Península) and onwards to Europe.
% Institutional Reception Centers (No Private Rental Access): During their temporary stay in the islands, they are housed in state-run emergency facilities (such as Las Raíces in La Laguna or Canarias 50 in Las Palmas) managed by NGOs (Cruz Roja, CEAR, Accem).
% Lack of Documentation: Without Spanish DNI/NIE, valid work permits, and employment contracts (nóminas), it is virtually impossible to sign private residential rental leases or secure mortgages.
% Unaccompanied Minors (Menores Tutelados): Those who remain under the custody of the Gobierno de Canarias are housed in public youth shelters and emergency care homes, not in the private real estate market.
% Urban Takeaway: Conflating the humanitarian crisis in emergency reception centers with the residential housing shortage is a common misconception. The housing shortage is driven by permanent residents, household unbundling, and planning paralysis.
% 3. The Pope's Visit
% Status: Factual in its humanitarian focus, but unrelated to local housing policy.
% The Context:
% Pope Francis explicitly announced his desire to visit the Canary Islands (reiterated in meetings with Canarian institutions) to spotlight the humanitarian drama of the Atlantic route, comparing it to his iconic 2013 visit to Lampedusa.
% His message focused strictly on human dignity, preventing the Atlantic from becoming a "cemetery without tombstones," and commending the humanitarian solidarity of the Canarian people.
% It was a moral and pastoral mission focused on sea rescues and human rights, not an urban-planning or housing-market critique.
% Summary: Demographic Drivers vs. Housing Impact in Santa Cruz de Tenerife
% Demographic / Visitor Flow  Volume  Nature of Stay  Impact on Santa Cruz Housing Stock
% Latin American & EU Migration  High & growing (Record foreign registration + Spanish descendants)  Permanent settlement (family homes, long-term employment)  Direct & High: Drives primary residential rental & purchase demand.
% Cruise Passengers  ~972.000 / year  Day-trippers (4–8 hours)  Zero: Sleep on ships; only impact street traffic & retail.
% Air Tourists (Canary total)  ~6.5M / year  Resorts (South/North) (Only ~1,600 stay in S/C hotels daily)  Marginal in S/C (except localized vivienda vacacional pressure).
% Irregular Sea Migration (Cayucos)  Humanitarian emergency  Transit (State reception facilities ➔ mainland transfers)  Negligible: Does not enter the private residential rental market.
% Internal Demographic Shift  Flat total population  Household Unbundling (more 1-2 person homes)  Very High: Same population requires thousands more dwellings.

% ==============================================================================
% ANÁLISIS DETALLADO: DINÁMICA DE FLUJOS MIGRATORIOS Y SALDO EXTERIOR (INE / ISTAC)
% ==============================================================================
% 1. Panel Superior: Entradas vs. Salidas Exteriores (Flujos Transatlánticos)
% - La crisis financiera española (2008–2014):
%   Las entradas desde el extranjero cayeron a su mínimo (~10.300 personas/año en 2009–2010).
%   Las salidas de residentes hacia el extranjero se duplicaron hasta alcanzar un máximo de 8.867 salidas en 2015 (salida de jóvenes y retorno temporal de extranjeros durante la recesión).
% - La inversión y ola migratoria de América Latina (2015–actualidad):
%   La profundización de la crisis institucional y económica en Venezuela, Colombia y Cuba provocó un giro histórico: las llegadas procedentes del exterior crecieron vertiginosamente de 14.199 (2015) a más de 32.800 personas/año (2022–2024).
%   Por el contrario, la emigración exterior se redujo notablemente, abriendo una brecha neta masiva a favor de las entradas.
%
% 2. Panel Inferior: Saldo Migratorio Neto Municipal (Santa Cruz de Tenerife)
% A escala estrictamente municipal (término de Santa Cruz de Tenerife, código 38038), los datos del ISTAC e INE muestran dos regímenes diametralmente opuestos:
% - Fase expulsora / pérdida neta (2009–2015):
%   Saldos fuertemente negativos: -2.168 hab. en 2010, -945 hab. en 2011 y -1.618 hab. en 2014. Santa Cruz perdía población por el freno inmigratorio y la salida hacia el exterior u otros municipios.
% - Fase de inversión y presión residencial sostenida (2016–2024):
%   A partir de 2016, el saldo municipal se vuelve intensamente positivo: +2.014 (2017), +2.545 (2018), +2.457 (2019), +2.084 (2022), +2.618 (2023) y +2.082 (2024).
%   Este superávit neto anual continuo de más de 2.000 nuevos residentes netos cada año en la capital explica directamente por qué la demanda habitacional permanente presiona el mercado de compra y alquiler, a pesar de que el padrón global total fluctúe de forma más moderada por el envejecimiento vegetativo de la población autóctona.
%
% 3. El trasfondo histórico: De "la Octava Isla" al retorno generacional
% - Flujo histórico (siglos XIX a 1980): Canarias fue durante más de un siglo una sociedad de emigración masiva hacia América (primero Cuba en los años 20 y luego Venezuela en los 50–70, popularizada como "la octava isla").
% - Flujo contemporáneo invertido: La inestabilidad en el continente americano trajo de vuelta a emigrantes retornados, a sus hijos y nietos (acogiéndose a la Ley de Memoria Democrática), así como a nuevos contingentes latinoamericanos.
% - Esto explica que el registro de extranjeros censados en Santa Cruz pasara de 4.473 en el año 2000 a 16.776 en 2025 (casi el cuádruple), a los que se suman decenas de miles que acceden directamente con pasaporte español por descendencia.
% ==============================================================================

Demographically, pressure on the housing stock also does not stem from humanitarian arrivals along the Atlantic maritime route—a transit phenomenon managed through institutional reception centers with zero entry into the private rental market.

Instead, residential housing demand is driven by the structural reversal of the historical transatlantic migration flow with the Americas:
Santa Cruz de Tenerife has transformed from an emigration society into a major recipient of permanent residents (particularly returnees and Latin American citizens with family roots or dual nationality).

The register of foreign nationals in Santa Cruz de Tenerife surged from 4,473 residents in 2000 to a record 16,776 in 2025, not including the tens of thousands who arrive holding Spanish citizenship by descent.

\begin{figure}[htbp]
  \centering
  \includegraphics[width=0.95\linewidth]{figures/migration_trend.pdf}
  \caption{Evolution of External Migration Flows and Net Migration Balance (2008--2024). Source: Author's elaboration based on INE (EMCR) and ISTAC (EVR) data.}
  \label{fig:migracion_trend}
\end{figure}

As reflected in official INE and ISTAC time series (Figure~\ref{fig:migracion_trend}), following the 2008--2014 recession cycle—characterized by declining arrivals and negative net balances as severe as $-2,168$ residents in the capital—external flows underwent a radical reversal from 2016 onwards. Foreign arrivals currently exceed 30,000 people annually across the province, consolidating an uninterrupted positive net migration balance in the municipality of Santa Cruz de Tenerife of over 2,000 permanent residents each year ($+2,618$ in 2023). This sustained influx of long-term residents exerts direct and immediate demand on the residential housing market, in an urban environment constrained by the lack of developable land.

\subsection{Families, New Households, and Housing Requirements}

% The Three Structural Shifts Driving Housing Demand
% The Math of Household Unbundling (Atomización Familiar):

% In the 1960s–70s (~3.4–3.8 people/household): To house Santa Cruz’s ~210,000 residents, the city needed around 61,000 dwellings.
% Today (~2.2 people/household, according to ISTAC): Housing those same ~210,000 residents requires more than 95,000 dwellings.
% Result: Even without population growth, demographic fragmentation alone generates an objective deficit of over 34,000 additional housing units.
% Hyper-Motorization and the "Garage Deficit":

% Tenerife has one of the highest motorization rates in Europe, with over 830–840 vehicles per 1,000 inhabitants (ISTAC). Nearly every adult in a working household owns or requires a vehicle.
% The Legacy Problem: Santa Cruz’s 20th-century housing stock (neighborhoods like Ofra, Los Gladiolos, La Salud, or García Escámez built between 1960 and 1980) was constructed with zero underground garages, assuming street-level parking for a largely carless working-class population.
% Today’s Requirement: Modern households consider at least one (often two) private parking spaces non-negotiable. Because older consolidated buildings cannot retrofit underground basements into steep volcanic topography, new demand concentrates exclusively on modern developments.
% Evolving Standards of Living Space (Superficie y Confort):

% Mid-century social housing (viviendas de absorción) crammed three bedrooms and a bathroom into 50–65 m².
% Contemporary emancipated households demand larger living spaces, dedicated work/study rooms for teleworking, storage rooms (trasteros), accessibility (elevators), and private terraces or balconies.
% The Territorial Consequence: Because the paralysis of the PGO prevents Santa Cruz from developing new residential sectors with modern standards and garages, younger generations are priced out or forced to migrate to peripheral municipalities (La Laguna, El Rosario, Candelaria), creating the daily traffic congestion on the TF-5 and TF-1 highways (las colas de la autopista).

% La necesidad de vivienda se ve agravada por el encogimiento del tamaño de los hogares: en 1960 cada vivienda ocupaba, en promedio, 3,4 personas, mientras que en la actualidad el promedio se sitúa en 2,2 personas, lo que exige disponer de más unidades para acoger a una población constante

Pressure on the real estate market of Santa Cruz de Tenerife is fundamentally amplified by three qualitative structural transformations within the urban fabric:

\begin{enumerate}
  \item \textbf{Household atomization and shrinking family size}: While in 1960 each home housed an average of 3.4 persons, the average household size in the municipality today stands at approximately 2.2 persons (ISTAC). This demographic fragmentation means that, to house a virtually stationary population (around 210,000 inhabitants), the city today requires more than 34,000 additional residential units compared to four decades ago, driven by the surge in single-person households and young adult emancipation.

  \item \textbf{Hyper-motorization and critical parking deficit}: With an island ratio exceeding 830 vehicles per 1,000 inhabitants, Tenerife has one of the highest motorization rates in Europe. Much of the consolidated housing stock built in the second half of the 20th century (working-class neighborhoods from the 1960s and 1970s such as Ofra, Los Gladiolos, or La Salud) was erected with zero underground parking. For contemporary households, the mandatory provision of one or two private parking spaces per dwelling is an unavoidable requirement that the older consolidated fabric cannot resolve, redirecting demand toward modern developments with basement garages.

  \item \textbf{Evolving standards of living space and habitability}: Contemporary spatial demands have rendered obsolete mid-century minimum social housing typologies (50 to 65 $m^2$ apartments designed for large families). Today's demand requires larger room dimensions, multi-purpose spaces for remote work, storage rooms (\emph{trasteros}), and universal accessibility.
\end{enumerate}

\begin{figure}[htbp]
  \centering
  \includegraphics[width=0.95\linewidth]{figures/income_trend.pdf}
  \caption{Evolution of Average Household Income in Santa Cruz de Tenerife (2011--2021). Source: Author's elaboration based on ISTAC and INE data.}
  \label{fig:renta}
\end{figure}

With municipal planning judicially blocked and expansion sectors paralyzed, the city has been unable to supply modern housing equipped with private garages. This bottleneck prices out newly formed households and pushes them toward the metropolitan periphery, worsening daily highway congestion along the TF-5 and TF-1 arterial corridors.
